% !TeX program = xelatex
\documentclass[tikz, margin = 5mm]{standalone}

\usepackage{xfrac}
\usepackage[MnSymbol]{mathspec}
\setmainfont[Mapping=tex-text,Ligatures={NoRequired,NoCommon,NoContextual}]{Comic Sans MS}
\setmathfont(Digits,Latin,Greek)[Numbers={Lining,Proportional}]{Comic Sans MS}
\usetikzlibrary{shapes.geometric,decorations,decorations.pathmorphing}
\usetikzlibrary{intersections}
\usetikzlibrary{patterns}
\usetikzlibrary{calc}
\usetikzlibrary{shapes.misc}
\usetikzlibrary{spy}

\pgfdeclarelayer{bg}
\pgfdeclarelayer{bbg}
\pgfsetlayers{bbg, bg, main}

\definecolor{c1}{HTML}{ac0000}
\definecolor{c2}{HTML}{00a0d5}

\newcommand{\abs}[1]{{\left|{#1}\right|}}

\tikzset{
	point/.style n args = {1}{circle, draw = black, inner sep = #1pt, fill = black},
	empty point/.style = {circle, draw = black, inner sep = 2pt, fill = white},
	cross/.style = {cross out, draw, minimum size = 2 * (#1 - \pgflinewidth), inner sep = 0pt, outer sep = 0pt},
	xkcd/.style n args = {1}{decorate, decoration = {random steps, pre length = #1mm, post length = #1mm, segment length = 0.6mm, amplitude = 0.2pt}},
	xkcd/.default = {4},
	point/.default = {2}
}

\begin{document}
	\begin{tikzpicture}
		\fill[xkcd, c2, fill opacity = 0.2] (-0.2,-0.25) rectangle (2.2,0.25);
		\fill[xkcd, c2, fill opacity = 0.2] (3.8,-0.25) rectangle (4.2,0.25);
		\fill[xkcd, c2, fill opacity = 0.2] (5.8,-0.25) rectangle (6.2,0.25);
		\draw[xkcd, thick] (-1,0) -- (7,0);
		\node[point, c1, label = {below}:{$4$}] (4) at (4,0) {};
		\node[point, c1, label = {below}:{$6$}] (6) at (6,0) {};
		\draw[xkcd = {0.1}] (0,-0.2) -- (0,0.2) node[pos=0, below] {$0$};
		\draw[xkcd = {0.1}] (2,-0.2) -- (2,0.2) node[pos=0, below] {$2$};
		\fill[xkcd, c1, fill opacity = 0.2] (0,-0.2) rectangle (2,0.2);
		\draw[xkcd = {1}, thick, ->] (0,0.5) -- (0.8,0.5);
		\draw[xkcd = {1}, thick, <-] (1.2,0.5) -- (2,0.5);
		\draw[xkcd = {1}, thick, ->] (3.4,0.5) -- (3.9,0.5);
		\draw[xkcd = {1}, thick, <-] (4.1,0.5) -- (4.6,0.5);
		\draw[xkcd = {1}, thick, ->] (5.4,0.5) -- (5.9,0.5);
		\draw[xkcd = {1}, thick, <-] (6.1,0.5) -- (6.6,0.5);
	\end{tikzpicture}
\end{document}